\documentclass[11pt]{article}

\usepackage[margin=1in]{geometry}
\usepackage{booktabs}
\usepackage{amsmath}
\usepackage{hyperref}

\title{11-768 HW2: Evaluating a Data-Visualization Agent}
\author{Your Name}
\date{}

\begin{document}
\maketitle

\section{Part 1: Examine Seed Agent Runs}

\subsection{Comparison of Ground-Truth and Validator Labels}
We evaluated the baseline validator (\texttt{validator/baseline.py}) across all 110 seed runs in \texttt{artifacts/} against the human ground-truth annotations in \texttt{seed\_labels.json}. Table~\ref{tab:part1-baseline} summarizes the per-family Matthews Correlation Coefficient (MCC), confusion counts, and positive/negative support for each error category.

\begin{table}[h]
\centering
\caption{Baseline validator performance on the 110 seed runs (\texttt{output/baseline\_predictions.json} vs.\ \texttt{seed\_labels.json}).}
\label{tab:part1-baseline}
\begin{tabular}{lrrrrrr}
\toprule
Error Family & MCC & TP & TN & FP & FN & Support ($+$/$-$) \\
\midrule
\texttt{execution\_failure} & $-0.0022$ & 4 & 65 & 33 & 8 & 12 / 98 \\
\texttt{wrong\_data}        & $0.0000$  & 0 & 73 & 0  & 25 & 25 / 73 \\
\texttt{wrong\_chart}       & $0.1096$  & 3 & 27 & 0  & 68 & 71 / 27 \\
\texttt{hard\_to\_read}     & $0.0000$  & 0 & 59 & 0  & 39 & 39 / 59 \\
\midrule
\textbf{Macro-MCC}          & \textbf{0.0268} & --- & --- & --- & --- & --- \\
\bottomrule
\end{tabular}
\end{table}

Overall agreement is near zero ($\text{macro-MCC} = 0.0268$), and disagreements concentrate in two sharply distinct regimes:
\begin{enumerate}
    \item \textbf{Bidirectional failure and cascading short-circuiting on \texttt{execution\_failure}:} Out of 12 true \texttt{execution\_failure} runs, the baseline detects only 4 (8 false negatives), while falsely flagging 33 valid runs as \texttt{execution\_failure} (33 false positives), yielding a slightly negative MCC ($-0.0022$). Because \texttt{validate()} returns immediately whenever \texttt{judge\_execution()} is non-empty, those 33 false positives also bypass \texttt{judge\_data\_and\_chart()} and \texttt{judge\_readability()}, automatically causing 3 \texttt{wrong\_data}, 14 \texttt{wrong\_chart}, and 8 \texttt{hard\_to\_read} ground-truth errors to be missed.
    \item \textbf{Extreme false-negative skew on content, design, and readability:} Even among the 65 valid runs that pass \texttt{judge\_execution()} and reach the remaining three checks, the baseline exhibits near-total recall collapse with zero false positives: it catches $0/22$ \texttt{wrong\_data} errors, $3/57$ \texttt{wrong\_chart} errors (runs \texttt{059}, \texttt{087}, and \texttt{173}), and $0/31$ \texttt{hard\_to\_read} errors, driving \texttt{wrong\_data} and \texttt{hard\_to\_read} MCCs to exactly $0.0000$.
\end{enumerate}

\subsection{Validator Error Type 1: Unreliable VLM \texttt{execution\_failure} Judgment and Cascading Short-Circuiting}
The baseline delegates \texttt{execution\_failure} entirely to the VLM via a negative-polarity question---\textit{``Did the agent crash, exhaust its turns, or give up without producing a valid figure?''}---and classifies a run as \texttt{execution\_failure} whenever the response begins with \texttt{"YES"} (\texttt{answer.upper().startswith("YES")}). In addition, \texttt{\_ask()} never inspects local disk state: when \texttt{run.figure} is \texttt{None}, it silently omits the image attachment without informing the model. This design causes both massive false positives and missed true failures:

\begin{itemize}
    \item \textbf{False-positive polarity inversion (33 of 33 FPs):} In every single one of the 33 false-positive \texttt{execution\_failure} predictions, the model's own generated \texttt{evidence} explicitly states that the agent \emph{succeeded} and produced a valid figure, yet the model opened its reply with \texttt{"YES"} (answering as if asked whether the task succeeded). Because \texttt{baseline.py} only checks the leading \texttt{"YES"} prefix and short-circuits, all subsequent error checks are skipped.
    \begin{itemize}
        \item \textbf{Run \texttt{004}:} In \texttt{artifacts/004/result.json}, \texttt{visualize\_functions.py} exited with \texttt{returncode: 0}, produced a $137{,}276$-byte 2$\times$2 \texttt{figure.png}, and exited with \texttt{exit\_status: "Submitted"}. Ground truth labels \texttt{004} as \texttt{wrong\_data} (panel 4 computes \texttt{y4 = 15 ** x4} without cubing) and \texttt{wrong\_chart} (panel 4 leaves the x-axis linear instead of logarithmic). However, the baseline returned \texttt{execution\_failure} with evidence: \textit{``The agent successfully completed the task... resulting figure.png file was created and verified to be valid... There were no crashes, turn exhaustion, or premature abandonment.''}, thereby missing both real errors.
        \item \textbf{Run \texttt{022}:} The agent wrote and ran \texttt{create\_figure.py}, producing a valid $42{,}528$-byte 1$\times$2 linear/log subplot \texttt{figure.png} and submitting cleanly. Yet the baseline again emitted \texttt{execution\_failure} while its evidence read: \textit{``The agent successfully completed the task. It created a Python script using matplotlib to generate a figure with two side-by-side subplots...''}.
    \end{itemize}
    \item \textbf{Missing-artifact blindness (8 of 12 FNs):} Across all 110 seed runs, all 12 ground-truth \texttt{execution\_failure} runs have no \texttt{figure.png} on disk (\texttt{run.figure is None}), whereas all 98 non-\texttt{execution\_failure} runs have a valid Pillow-decodable \texttt{figure.png}. Because \texttt{baseline.py} does not check \texttt{run.figure is None} or tell the VLM that the image is absent, the VLM frequently hallucinates success from trajectory text alone.
    \begin{itemize}
        \item \textbf{Run \texttt{010}:} The agent looped for 23 bash commands (repeatedly activating a virtualenv and failing to produce a figure) until hitting \texttt{exit\_status: "LimitsExceeded"}, leaving no \texttt{figure.png}. The baseline predicted \texttt{[]} (no errors).
        \item \textbf{Run \texttt{096}:} The agent's \texttt{plot.py} saved to the wrong filename (\texttt{'create figure.png'}) rather than \texttt{figure.png}, so \texttt{run.figure} is \texttt{None} at evaluation time. Instead of reporting \texttt{execution\_failure}, the baseline missed it and paradoxically predicted \texttt{wrong\_chart}.
    \end{itemize}
\end{itemize}

\subsection{Validator Error Type 2: Systematic False Negatives from Raw Trajectory Noise and Unstructured Single-Shot Prompts}
Among the 65 valid runs that actually reached \texttt{judge\_data\_and\_chart()} and \texttt{judge\_readability()}, the baseline missed 107 out of 110 ground-truth error instances (22/22 \texttt{wrong\_data}, 54/57 \texttt{wrong\_chart}, and 31/31 \texttt{hard\_to\_read}). Two structural flaws in \texttt{validator/baseline.py} explain this conservatism:
\begin{enumerate}
    \item \textbf{Raw, noisy trajectory serialization:} \texttt{\_ask()} dumps \texttt{str(run.messages)} verbatim (median $16{,}547$ characters, up to $79{,}883$ characters), burying the final executed plotting script inside escaped JSON strings, shell tracebacks, tool-call IDs, and superseded earlier scripts that often look correct in parts where the final \texttt{plot.py} regressed.
    \item \textbf{Vague single-sentence \texttt{YES/NO} prompts without rubric or reasoning-before-verdict:} Each check asks a single broad question (e.g., \textit{``Does the plotted data differ from what was requested in any way?''} or \textit{``Is the rendered figure difficult or impossible to read?''}) with no definition of what counts under \texttt{wrong\_data}, \texttt{wrong\_chart}, or \texttt{hard\_to\_read}, and forces the model to output \texttt{YES} or \texttt{NO} as its very first token before inspecting individual requirements. When confronted with a dense multi-thousand-token blob and an underspecified question, the 30B VLM overwhelmingly defaults to \texttt{NO}.
\end{enumerate}

Four representative runs that passed \texttt{judge\_execution()} yet received \texttt{[]} (all errors missed) illustrate this failure across all three downstream families:
\begin{itemize}
    \item \textbf{Run \texttt{008} (GT: \texttt{wrong\_data}, \texttt{hard\_to\_read}; Pred: \texttt{[]}):} In the final executed \texttt{plot.py} (message 14), the agent explicitly filtered out \texttt{orange} from sources via \texttt{.query("source != 'orange'")}, dropping the entire \texttt{orange} source node present in \texttt{data.csv} (even though an earlier unused draft \texttt{create\_sankey.py} in the same trajectory had included all labels). It also rendered every node value label in near-white \texttt{color="\#f4f4f4"} on a white canvas, making the labels virtually invisible in \texttt{figure.png}. The baseline missed both defects.
    \item \textbf{Run \texttt{009} (GT: \texttt{wrong\_chart}, \texttt{hard\_to\_read}; Pred: \texttt{[]}):} The instructions explicitly require \textit{``Rotate the x tick labels by 30 degrees''}. However, the final \texttt{plot.py} sets \texttt{mdates.MonthLocator(interval=1)} with \texttt{rotation=0}, producing an illegible horizontal black smear of overlapping date labels along the bottom axis of \texttt{figure.png}. Both the code violation and visual overlap were missed.
    \item \textbf{Run \texttt{013} (GT: \texttt{wrong\_chart}; Pred: \texttt{[]}):} The requirements explicitly state \textit{``do NOT use the jet, rainbow or hsv colormaps''} and \textit{``Set the x axis limits to 0..20''}. The final \texttt{plot.py} literally calls \texttt{plt.colormaps["rainbow"]} and \texttt{ax.set\_xlim(-5, 25)}, both plainly visible in the script and the 3D plot, yet the baseline answered \texttt{NO}.
    \item \textbf{Run \texttt{016} (GT: \texttt{wrong\_data}, \texttt{wrong\_chart}; Pred: \texttt{[]}):} The task requires dropping browser versions below $0.5\%$ share and showing a browser legend. An earlier failing draft (\texttt{create\_chart.py}) contained \texttt{df[df['Data'] >= 0.5]}, whereas the final executed \texttt{plot.py} omitted the filter (plotting IE~9.0 at $0.44\%$, UC Browser~7.0 at $0.36\%$, and Chrome~65.0 at $0.25\%$) and omitted \texttt{ax.legend()}. Because the raw trajectory mixed the superseded draft with the final script under a vague prompt, the baseline failed to catch either error.
\end{itemize}

\subsection{Proposed Validator Changes}
To address both failure modes in Part~2 without hard-coding task IDs or replacing the fixed validator model, we propose two targeted, generalizable changes in \texttt{validator/solution.py}:

\begin{enumerate}
    \item \textbf{Change 1 (Deterministic \texttt{execution\_failure} Verification):} Replace the polarity-prone VLM prompt in \texttt{judge\_execution(run)} with deterministic local Python and Pillow checks. Specifically, verify that \texttt{run.figure} exists, has non-zero byte size, decodes cleanly as an image via \texttt{PIL.Image}, has non-degenerate pixel dimensions, and contains non-uniform pixel content (not a blank/single-color canvas). This eliminates both missing-figure false negatives (e.g., runs \texttt{010}, \texttt{096}) and false-positive short-circuiting on valid runs (e.g., runs \texttt{004}, \texttt{022}) at zero VLM cost.
    \item \textbf{Change 2 (Structured Trajectory Extraction + Rubric-Guided Analyze-Then-Verdict Prompting):}
    \begin{itemize}
        \item \emph{Structured input curation:} Parse \texttt{run.messages} to separate the final executed plotting script (and its terminal output/exit status) from superseded earlier drafts, and attach deterministic image metadata (canvas pixel size and DPI from Pillow) alongside bounded input-data previews. Place the shared run context before per-category audit instructions (improving both signal-to-noise ratio and vLLM prefix cache reuse).
        \item \emph{Rubric-guided JSON reasoning:} Replace the bare leading-\texttt{YES/NO} questions for \texttt{wrong\_data}, \texttt{wrong\_chart}, and \texttt{hard\_to\_read} with explicit operational rubrics adapted from \texttt{ASSIGNMENT.md} (checking values/filtering/ordering against the final script, checking chart type/axes/limits/colormap/legend/labels against every requirement, and inspecting the rendered image for clipping, overlap, occlusion, contrast, and indistinguishability). Require the VLM to first list requirement-by-requirement observations from both the final script and image, and only then emit a structured JSON verdict (\texttt{has\_error: bool} and grounded \texttt{evidence}).
    \end{itemize}
\end{enumerate}

\section{Part 2: Improve Your Validator}

\subsection{Implemented Changes}
Describe the two changes implemented from Part 1. Distinguish deterministic
checks from model-based checks, where applicable, and explain how each change
uses the task, trajectory, input data, or rendered figure.

\subsection{Seed-Set Performance vs.\ Baseline}
Report the modified validator's performance on the seed runs alongside the
baseline using all four per-family MCC scores and macro-MCC.

\subsection{Qualitative Analysis of Targeted Error Types}
For each of the two error types surfaced in Part 1, discuss whether the
corresponding change actually mitigated it, with reference to specific runs.
Also describe any regressions.

\section{Part 3: Design New Evaluation Tasks}

\subsection{Anticipated Generalization Weaknesses}
Identify two classes of visualization task that your current validator may
have trouble generalizing to, and explain why.

\subsection{Task Design}
Describe the five visualization tasks, including at least two from each
weakness class. For each task, describe its inputs and data source, expected
behavior, and why it is a useful stress test for the validator.

\subsection{Agent Runs and Labeling}
Describe the agents run on the new tasks, the runs generated, and how you
assigned human labels using the four fixed error families. Explain any
ambiguous labeling decisions.

\subsection{Validator Performance on New Tasks}
Report all four per-family MCC scores and macro-MCC on the new agent runs and qualitatively
analyze important successes, failures, and uncertain labels.

\subsection{Harbor Packaging}
State that every authored task scores 1.0 against \texttt{-a oracle}; you may
also report the recommended \texttt{-a nop} runs. Describe both plausible-but-wrong
mutant solutions and report each reward. For the mutant scoring 0, explain what
it gets wrong and name the assertion that rejects it. For the mutant scoring
1.0, explain what it gets wrong and why the verifier accepts it. Then, for each
of the four error families, say whether a deterministic Harbor verifier for your
task could decide it, and why or why not. Finally: your validator is graded on
private runs for tasks you never authored. Explain why a per-task verifier
cannot serve as that grading mechanism.

\section{Part 4: Iterate on Your Design}

\subsection{Additional Improvement Motivated by Part 3}
Describe at least one change to the validator motivated by the Part 3 results.
Identify the observed failure that motivated it and explain why you expected
the change to help.

\subsection{Final Evaluation on Seed and New Tasks}
Evaluate the final validator on both the seed runs and self-authored runs.
Report all four per-family MCC scores and macro-MCC before and after the change on each subset,
including any regressions.

\subsection{Qualitative Analysis of Self-Authored Tasks}
Discuss whether the change improved performance on your self-authored tasks,
with reference to specific runs.

\subsection{Further Changes}
Report any additional validator changes or tasks beyond the required ones.
For each, explain its motivation and how successful you judged it to be,
qualitatively and, for quantitative comparisons, using all four per-family MCC
scores and macro-MCC.

\end{document}
